\section{Discussion}
\label{sec:discussion}

\subsection{What the Loop Buys, What the Ladder Prices, and What Neither Can}
The verified loop replaces confidence in the generator with an audited
pipeline, and the evaluation localizes what each stage
contributes. The ladder adds the accounting: which rungs of intelligence
carry measurable weight. The RQ4 results cut both ways -- the
grounding rung is load-bearing, while the top rung's routing benefit
was not resolvable above the deterministic floor -- and the rung's
provider refused assessment precisely on the highest-stakes sessions --
so the LLM came out of that path. The loop makes spent intelligence
safe; the ladder keeps it cheap. Neither makes the provider problem
disappear, and the hard-subset census adds a wrinkle: the free-tier
endpoint intermittently refuses security-shaped payloads outright
(HTTP 403 on attacker-tooling states, benign states served). For a
deployment whose critical moments are attack-shaped, provider refusal
is not only a cost risk but a selective availability risk concentrated
exactly where assessment matters most. Every
generative claim in this paper is downstream of a free-tier inference
API: the safety-alignment timeouts that shaped the persona stratum, the
catalog churn the watcher caught, and the quota windows that gated the
study schedule are all the same structural fact. We treat the provider as
an unreliable component and the system degrades safely when it fails, but
a practitioner sizing a deployment should budget for generation being
\emph{eventually unavailable} in exactly the way the fallback path
assumes -- and should read the fold-down rule as the remedy: every
mechanism compiled down to a certifiable equivalent is one fewer
dependency on that provider.

\subsection{Deployment Guidance}
Five lessons from operating the system generalize. First, run the
watcher: provider catalogs drift, and a wired model disappearing silently
is a production incident for any LLM-dependent defense. Second, scope the
HITL gate with care: the gate exists for content classes whose
compromise value is high, and the audit log it produces doubles as the
governance record. Third, keep the deterministic floor maintained: the
fallback transcripts are load-bearing. The system is only as consistent
as its last vetted transcript, and treating those transcripts as second-class
content would hollow out the safety property that makes the LLM path
permissible. Fourth, budget for provider instability as a first-class
operating condition: during this study the watcher logged one wired model
disappearing from the catalog and one quota-exhaustion window, and the
fallback path absorbed both without serving unverified content. An
LLM-dependent defense sized without that budget will be sized wrong.
Fifth, audit the ladder on a schedule: re-price each intelligent rung
against its cheapest certifiable competitor as the deterministic
machinery matures. In this study the audit already paid out once -- the
kind-aware routing rung was removed after deterministic fixes reproduced
its accuracy offline, leaving the serving path with one fewer model,
provider, and failure mode -- and in the released system the audit is
not a manual exercise: the judge-panel audit and audited stability
reseed automatically whenever a new campaign run lands (Appendix~\ref{app:audit}).

\subsection{Cost and Practicality}
A defender considering this architecture is entitled to ask what it costs.
The serving host is a commodity box (2~vCPU, ${\sim}4$~GB RAM) shared with
the honeypot listeners; profiling adds milliseconds per session, retrieval
over the curated corpus adds tens of milliseconds, and the dominant cost
is provider inference, bounded per draft by output caps and low
temperature. The mean session in the live campaign completed in
${\sim}44$~s end to end, dominated by the generative turns, and the
deterministic strata served in fractions of a second. The economically
interesting number is the fallback rate: when the provider refused or
timed out, the loop degraded to vetted transcripts rather than
open-ended retries, so the cost ceiling per session is two generative
attempts plus one deterministic serve. On a free-tier allowance the full
218-query campaign completes within a daily quota window, which is the
difference between a system a lab can operate and one it cannot.

\subsection{Generalizing the Ladder and the Loop}
Nothing in the loop is specific to honeypots, and neither is the ladder.
The shape, generate under
grounded context, verify the draft against a persisted state structure,
serve only on pass, fall back to vetted content on failure, applies
wherever an LLM writes output whose downstream consumer is adversarial or
safety-critical: SOC copilot drafting incident narratives, retrieval
chatbots answering from a curated corpus, agent pipelines writing into
structured stores. The fiction graph generalizes to any \emph{state
bound} the LLM must respect; the verifier to any checkable draft
property; the deterministic fallback to any canned-but-correct content.
The certifiability axis is the portable part of the ladder: any
mechanism in any such pipeline can be placed on it, priced by marginal
effect, and audited for fold-down as certifiable alternatives mature.
The evaluation here does not establish that the
pattern transfers, but the
architecture section is written so that a practitioner in another domain
can lift the loop wholesale, and the ablation battery (A1--A5) plus the
RQ4 pricing protocol are the template for testing whether each piece
they keep contributes.

\subsection{Ethics}
The system deploys deception against real attackers, which obligates
care. All credential material shown to attackers is registered canary
secrets: nothing real is ever exposed, and any canary appearing in
outbound traffic is attributable to bait interaction by construction. The
HITL gate keeps a human decision between the system and the highest-stakes
content classes, and the append-only audit log records every gate
decision. The testbed is isolated from production networks, and the
deception targets only actors actively attacking it. We release scenario
sets and scoring harnesses so that the community can audit the
evaluation, not just the code.

\subsection{Limitations}
The deployment observes tailnet traffic only; the ambient-internet axis
(scanners, botnets, opportunistic ransomware) is absent by design choice
and limits ecological claims. The consensus study's three evaluator
personas share one judge model family; multi-model judging failed on the
provider at study time (a listed model returned 404 on real calls despite
catalog presence), a gap we flag for the watcher's future checks. The
$0.750$ agreement figure should therefore be read as within-model
consistency -- how stably one model's prompted personas concur -- rather
than inter-annotator agreement across independent judges. The shell-scope listener constrains web-shaped deception scenarios, as the
weakest consensus scenario illustrates. And the persona-realism stratum,
while stable run-to-run, remains the stratum where judge variance lives;
its scores should be read with the MAD in \S\ref{sec:stability} in mind.
The RQ4 evidence carries its own weight limits: the real-log assessment
sample is 13 sessions, the closed-loop replay pairs number five
clean families, and the simulated attacker, however instrumented, is not
an adversary with a stake in the outcome. The routing fold-down rests on
200 held-out sessions at ceiling plus a full-corpus census locating the
informative stratum (six sessions, one rescued live, five lost to
provider refusals); it is stated at the benchmark's resolution rather
than as general subsumption. The deception-effect claims for the
top rung (bait exposure, the dwell null) are directional, and we price
them as such rather than as established effects.
